% Source preserved in full: originales/cuaderno/NOTA.md, §§1--7.
\section{Chronological memory, balance, and vacancy response}
\label{md:cron:seccion}

The following construction receives the symbolic capacity factor and studies
its subsequent readout operators. The internal ratio $3^6/10^3=729/1000$ and its
calendar belong to the preceding construction
(Sections~\ref{mono-ampl:relojes} and~\ref{vac-capacidad-trit}).
The APP--TRIT--TPK architecture, the regions, nonadic return, and the
record $K$ retain their prior causal position. The logarithms used
in the capacity readout operator describe that factor and do not select the
regions of $\pi,\varphi,e,\alpha$ anew.
The propositions are focal consequences concerning observation and memory;
their scope is distinguished from the full TPK state and from an attribution
of historical priority for general results on mechanical sequences or
information.

\subsection{Object, measure, and two readout operators}
\label{md:cron:lectores}

Let $\delta=\log_{10}(10/9)$ be the slope of the capacity calendar,
and let $\theta$ be a phase uniformly distributed on $[0,1)$.
The measure represents uncertainty about the readout origin, while the
mechanism remains determined. Define
\begin{equation}
 \begin{aligned}
 z_j(\theta)&=\lfloor(j+1)\delta+\theta\rfloor
             -\lfloor j\delta+\theta\rfloor,\\
 W_n(\theta)&=(z_0(\theta),\ldots,z_{n-1}(\theta)),\\
 B_n(\theta)&=\sum_{j=0}^{n-1}z_j(\theta).
 \end{aligned}
 \label{md:cron:dos-lectores}
\end{equation}
The block $W_n$ retains chronology, and $B_n$ retains balance.
For a known exact phase and exact rule, both objects are
deterministic; the entropies refer expressly to the set of possible
phases.

Telescoping gives
\begin{equation}
 B_n(\theta)=\lfloor n\delta+\theta\rfloor.
 \label{md:cron:telescopia}
\end{equation}
Thus $B_n$ takes only the values $\lfloor n\delta\rfloor$ and
$\lceil n\delta\rceil$. The points $\{-j\delta\}$, for $j=0,\ldots,n$,
divide the circle into $n+1$ phase intervals. Each interval carries a
distinct sequence of length $n$. Indeed, the cumulative prefixes
$\lfloor k\delta+\theta\rfloor$, for $k=1,\ldots,n$, are monotone in
$\theta$, and each boundary changes one of them; the sequence reconstructs
all those prefixes.

The endpoints have measure zero. A half-open convention resolves their
assignment and preserves the entropies.

\subsection{Increasing information and zero entropy rate}
\label{md:cron:crecimiento}

Let $\ell_0,\ldots,\ell_n$ be the lengths of the phase intervals, and let
\begin{equation}
 H_n=-\sum_{j=0}^{n}\ell_j\log_2\ell_j.
 \label{md:cron:entropia-bloque}
\end{equation}

\begin{proposition}
\label{md:cron:proposicion-tasa}
The following hold simultaneously:
\begin{equation}
 H_n\longrightarrow+\infty,\qquad
 \frac{H_n}{n}\longrightarrow0.
 \label{md:cron:dos-limites}
\end{equation}
\end{proposition}
\begin{proof}
The upper bound is $H_n\leq\log_2(n+1)$. If
$\ell_{\max}(n)=\max_j\ell_j$, then
$H_n\geq-\log_2\ell_{\max}(n)$.
The slope $\delta$ is irrational: if $\delta=p/q$ with $q>0$,
then $(10/9)^q=10^p$, contradicting the valuation at the prime $3$.
By irrationality, the orbit points are dense. The successive
partitions have maximum mesh tending to zero: for every $\epsilon>0$
there is a finite $\epsilon$-dense segment, and later segments only
subdivide its intervals. The lower bound therefore diverges. The
upper bound divided by $n$ tends to zero.
\end{proof}

The interval lengths are not equal in general. The proof
establishes the preceding bounds without asserting
$H_n=\log_2(n+1)$ or an exact logarithmic asymptotic law that would require
additional Diophantine conditions.
The topological entropy of the factor is also zero because
$\log(n+1)/n\to0$. Block information and its rate are different
quantities. Growth of the archive of a known trajectory does not,
by itself, amount to production of independent statistical information.

\subsection{Chronological information not determined by balance}
\label{md:cron:condicional}

\begin{corollary}
\label{md:cron:corolario-condicional}
For the same set of phases,
\begin{equation}
 H(W_n\mid B_n)=H_n-H(B_n)
 \geq H_n-1\longrightarrow+\infty.
 \label{md:cron:entropia-condicional}
\end{equation}
\end{corollary}
\begin{proof}
$B_n$ is a function of $W_n$, so the chain rule gives the equality.
Its support has two elements, hence $H(B_n)\leq1$ bit.
Apply Proposition~\ref{md:cron:proposicion-tasa}.
\end{proof}

The conditional entropy that diverges is the average over balances;
the same uncertainty is not assigned to every particular balance.
Nor is that loss attributed to the full enriched state: when the
retained memory recovers $W_n$, one has
\begin{equation}
 H(W_n\mid B_n,M_n)=0.
 \label{md:cron:restitucion-memoria}
\end{equation}
By cardinality, at least one of the two fibres of $B_n$ contains
$\lceil(n+1)/2\rceil$ sequences or more. This combinatorial bound is
distinct from the entropy bound.
The result specifies the relation between value and genealogy for this readout operator:
the two possible balances do not exhaust the compatible temporal histories.

\subsection{Centred current: bounded balance and persistent activity}
\label{md:cron:corriente}

The centred readout operator is expressed by
\begin{equation}
 J_j=z_j-\delta,\qquad
 \sum_{j=0}^{n-1}J_j=B_n-n\delta.
 \label{md:cron:corriente-definicion}
\end{equation}
Consequently,
$\left|\sum_{j=0}^{n-1}J_j\right|<1$ for every window.
At the same time,
\begin{equation}
 \lim_{n\to\infty}\frac1n\sum_{j=0}^{n-1}J_j=0,\qquad
 \lim_{n\to\infty}\frac1n\sum_{j=0}^{n-1}J_j^2
 =\delta(1-\delta)>0.
 \label{md:cron:actividad}
\end{equation}
\begin{proof}
One has $z_j^2=z_j$ and $B_n/n\to\delta$ by the discrepancy bound.
Expanding
$(z_j-\delta)^2=z_j(1-2\delta)+\delta^2$ and substituting the
average of $z_j$ gives the second identity.
The argument holds for every phase and requires no random averaging.
\end{proof}

Bounded cumulative balance, positive quadratic activity, and zero
entropy rate are therefore compatible. The identities preserve the
distinction between the centred readout operator and a representation of elementary
charge; nor do they establish energy production without a balance.

If an already established dimensional realization takes
$I_j=(u_Q/t_0)J_j$, one obtains
\begin{equation}
 I_{\mathrm{rms}}^2
 =\left(\frac{u_Q}{t_0}\right)^2\delta(1-\delta).
 \label{md:cron:eficaz}
\end{equation}
If each $I_j$ is maintained for an interval of duration $t_0$, this
discrete average coincides with the time average.
As a subsequent physical comparison, an ideal resistive load
$R_{\mathrm{res}}>0$ would have mean power
\begin{equation}
 \overline P_{\mathrm{res}}
 =R_{\mathrm{res}}\left(\frac{u_Q}{t_0}\right)^2
 \delta(1-\delta)>0.
 \label{md:cron:potencia-resistiva}
\end{equation}
This is a prediction of the model composed with that load. The resistance and
the device retain their own specifications; they do not follow from the
capacity sequence alone.

\subsection{Current spectrum and orientation}
\label{md:cron:espectro}

The observation function is
$f(\theta)=\mathbf1_{[1-\delta,1)}(\theta)-\delta$, with
$J_j=f(\theta+j\delta)$. Its Fourier coefficients, for $m\neq0$, are
\begin{equation}
 a_m=\frac{\exp(2\pi i m\delta)-1}{2\pi i m},\qquad
 |a_m|^2=\frac{\sin^2(\pi m\delta)}{\pi^2m^2},\qquad
 a_0=0.
 \label{md:cron:fourier}
\end{equation}
They are obtained by integrating the indicator function over its interval. Parseval's
identity gives
\begin{equation}
 \sum_{m\neq0}|a_m|^2=\delta(1-\delta),
 \label{md:cron:parseval}
\end{equation}
in agreement with quadratic activity. The temporal frequencies
are $m\delta$ modulo $1$. The stationary correlation, with uniform phase,
is determined by these intensities; they do not contain the initial phase.
A translation multiplies $a_m$ by a unit-modulus phase factor, and time
reversal exchanges opposite frequencies.

Spectral power alone recovers neither orientation nor the temporal
origin. A phase-sensitive readout or a marked reference supplies
the complementary information. The deduction preserves the distinction between
that general reference and the record $K$, as well as between symbolic rhythms
and a particular experimental realization.

\subsubsection{Cyclic reference and transports of equal spectral intensity}
\label{md:cron:k-ciclico}

The cyclic property of the record $K$ yields a finite consequence
independent of the preceding spectrum. In $V_K=\mathbb R^{12}$, let $S$ be the
shift $(Sx)_i=x_{i+1}$ and let
\begin{equation}
 O_K=(K,SK,\ldots,S^{11}K).
 \label{md:cron:matriz-ciclica}
\end{equation}
We use the invertibility of $O_K$ established by the cyclic reference
of the record. For any real operator $H$ commuting with $S$,
put $y_+=HK$ and $y_-=H^TK$.

\begin{proposition}
\label{md:cron:proposicion-intensidades}
The twelve Fourier intensities of $y_+$ and $y_-$ coincide.
However, $y_+=y_-$ if and only if $H=H^T$.
Each full vector response, together with $K$ and the orientation of $S$,
recovers its operator.
\end{proposition}
\begin{proof}
$H$ is circulant. Its Fourier multipliers are $\lambda_m$, and those
of $H^T$ are their conjugates, so the two responses have
intensities $|\lambda_m|^2|\widehat K_m|^2$.
If their vectors coincide, $(H-H^T)K=0$. Commutation implies that
$H-H^T$ also annihilates $S^jK$ for every $j$. These vectors form a basis;
therefore $H-H^T=0$. The converse is immediate.
Recovery is expressed by
\begin{equation}
 H=O_{y_+}O_K^{-1},\qquad
 H^T=O_{y_-}O_K^{-1}.
 \label{md:cron:restitucion-operador}
\end{equation}
\end{proof}

In particular, the shifts $S$ and $S^{-1}$ produce responses with
the same power spectrum but different vectors. The exact autocorrelations
of the record give
\begin{equation}
 \|SK-S^{-1}K\|^2
 =2\bigl(\|K\|^2-\langle K,S^2K\rangle\bigr)
 =2175744>0.
 \label{md:cron:separacion-desplazamientos}
\end{equation}
Rational inversion of $O_K$ recovers from the two responses the
coefficients $e_1$ and $e_{11}$, respectively. $K$ is used after
its generation as a marked reference. Spectral intensity or scalar
balance does not replace the vector response.
The corollary describes the twelve-component cyclic carrier;
it preserves the difference between $S$ and the full enriched dynamics,
as well as the additional operators of a physical charge conjugation.
For general $H$, transposition does not mean time reversal either;
for $S$, it does coincide with $S^{-1}$.

\subsection{Electronic operator link}
\label{md:cron:electron}

Incidence between the APP sum and product fibres over $D_9^3$
defines $C$. Its tritic vector
$m_{\mathrm{ph}}=(1,-1,0)^3$ satisfies
\begin{equation}
 Cm_{\mathrm{ph}}=C^Tm_{\mathrm{ph}}
 =-\frac12m_{\mathrm{ph}}.
 \label{md:cron:modo-tritico}
\end{equation}
The induced calendar traverses the regimes of the same tritic selector;
the incidence selects the principal plane with singular value $1/2$.
On that plane,
\begin{equation}
 P=\begin{pmatrix}1&0\\0&0\end{pmatrix},\qquad
 Q=\begin{pmatrix}1/4&\sqrt3/4\\\sqrt3/4&3/4\end{pmatrix},
 \qquad J=\frac4{\sqrt3}[P,Q],\qquad J^2=-I.
 \label{md:cron:plano-electronico}
\end{equation}
The incidence and the commutator norm are developed in
Section~\ref{vi:dep:prefactor}. An effective relation between
temporal organization and the electronic block is thereby recovered, preserving the selector,
its labels, and the incidence. The charge representation and the dimensional
realization retain their own operators: the link does not assign
opposite elementary charges directly to $z_j-\delta$.

\subsection{Retained information and energy composition}
\label{md:cron:informacion-energia}

Fibre recovery establishes
\begin{equation}
 H(X\mid q(X),M(X))=0
 \label{md:cron:inversa-enriquecida}
\end{equation}
when the enriched readout operator has an inverse; a bijective update
preserves $H(X)$. These properties are assembled in
\eqref{eq:ii-kb-aut-fibra} and~\eqref{eq:ii-kb-aut-reversible}.
The coexistence of increasing $H_n$ and zero rate is compatible with those
properties and with Shannon's identities.

Erasure work depends on the accessible information retained.
In the ideal isothermal regime of degenerate logical memories and in the
appropriate asymptotic limit, the cost per copy with side information is
\begin{equation}
 W_{\mathrm{borrado,ideal}}
 =k_BT\ln(2)\,H(X\mid M).
 \label{md:cron:borrado}
\end{equation}
If $X$ can be recovered from $M$, the ideal logical contribution may be zero.
Physical transport, readout, control, and the full cycle retain
their own balances; the equality $h_{\mathrm{top}}=0$ does not remove them.

The relevant composition preserves the sequence
\[
 \begin{gathered}
 \text{chronology}\longrightarrow
 \text{boundary record}\longrightarrow
 \text{current and state readout operator}\\
 \longrightarrow
 \text{retained memory}\longrightarrow
 \text{thermal balance}.
 \end{gathered}
\]
The preceding development constructs the first two readout operators and their
informational consequences. Its scope is distinguished from certification
of a complete physical device and leaves intact the constructive
genealogy of the constants.
